\chapter{Couche 10 — Les limites assumées}

Ce chapitre est le plus important du document. Une défense dont on ignore la
limite donne une fausse assurance, et une fausse assurance est pire qu'une
absence de défense : on arrête de surveiller.

\section*{Les risques résiduels, et pourquoi ils sont acceptés}

\begin{center}
\small
\begin{tabularx}{\textwidth}{@{}p{3.6cm}Xp{1.9cm}@{}}
\toprule
\textbf{Risque} & \textbf{Pourquoi il est accepté ici, et ce qu'il faudrait en production} &
\textbf{Gravité} \\
\midrule
\textbf{Certificat auto-signé} &
Le trafic est chiffré, mais l'identité du serveur n'est pas prouvée : un
attaquant capable de s'interposer pourrait se faire passer pour la plateforme.
Seule alerte que l'analyse dynamique signale encore. Correction prévue :
Let's Encrypt. & Moyenne \\
\addlinespace
\textbf{VncAuth} &
Défi-réponse DES, mot de passe tronqué à huit caractères par le protocole.
Acceptable parce que la machine est éphémère, que le mot de passe n'est visible
que cinq minutes et que le transport est chiffré. En production : un proxy
authentifié devant noVNC. & Moyenne \\
\addlinespace
\textbf{\cmd{docker.sock} monté dans la pile compose} &
En mode mono-hôte, le conteneur de l'application reçoit la socket Docker de
l'hôte --- donc un pouvoir équivalent à \cmd{root} sur la machine. Le fichier
\cmd{docker-compose.yml} le dit en toutes lettres : \emph{« À n'utiliser que
sur une machine de développement. »} Sur le parc déployé, ce montage
\textbf{disparaît} : le contrôleur pilote les workers par SSH. & Élevée
(développement seulement) \\
\addlinespace
\textbf{\cmd{root} dans le conteneur loué} &
C'est le service vendu, pas un défaut. Atténué par les capacités retirées,
\cmd{no-new-privileges} et les limites de ressources. Une évasion par faille du
noyau reste possible : seule une machine virtuelle créerait la vraie
frontière. & Acceptée \\
\addlinespace
\textbf{Un seul jeton pour tous les agents} &
Il ne distingue pas les nœuds entre eux ; le compromettre permet d'en déclarer
un faux, donc de recevoir des machines à héberger. Une authentification
mutuelle par certificat demanderait une autorité de certification. & Moyenne \\
\addlinespace
\textbf{La reprise ne restitue pas les fichiers} &
Elle recrée un conteneur, elle ne déplace pas un disque. Sans stockage partagé
entre nœuds, seul le service repart. C'est le premier élément des travaux à
moyen terme. & Acceptée \\
\addlinespace
\textbf{Les clés sur le disque du contrôleur} &
\cmd{.secret\_key} et \cmd{.vault\_key} sont en \cmd{0600}, mais sur la même
machine que la base. La séparation des deux clés est logique, pas physique. Un
gestionnaire externe les garderait hors de la machine. & Moyenne \\
\addlinespace
\textbf{Pas de journalisation centralisée} &
Les traces d'accès aux machines louées restent sur chaque nœud. Une
investigation après incident serait laborieuse, et un attaquant ayant pris un
nœud peut effacer ses traces. Travaux à moyen terme. & Moyenne \\
\addlinespace
\textbf{Pas de second facteur} &
Un mot de passe compromis suffit à prendre un compte. Travaux à moyen terme. &
Faible (contexte pédagogique) \\
\bottomrule
\end{tabularx}
\end{center}

\section*{Ce qui est hors de portée du projet}

Trois choses qu'aucune mesure du document ne traite, et qu'il serait malhonnête
de laisser croire couvertes :

\textbf{Un attaquant déjà administrateur du contrôleur} lit la base, les clés,
et pilote tous les nœuds. La défense est organisationnelle --- qui a accès à la
machine --- pas logicielle.

\textbf{Une vulnérabilité du noyau Linux} traverse la frontière des conteneurs,
qui partagent ce noyau avec l'hôte. Le durcissement réduit la surface, il ne
crée pas d'isolation matérielle.

\textbf{Un déni de service volumétrique} venu de l'extérieur relève de
l'hébergeur. Les limites par conteneur protègent les locataires les uns des
autres, pas la plateforme d'une attaque réseau.

\section*{Récapitulatif : menace, parade, preuve}

\begin{center}
\small
\begin{tabularx}{\textwidth}{@{}p{4.1cm}p{4.4cm}X@{}}
\toprule
\textbf{Menace} & \textbf{Parade} & \textbf{Preuve} \\
\midrule
Fuite de la base & PBKDF2 salé, Fernet &
Les mots de passe en base sont illisibles sans la clé du coffre. \\
Énumération des comptes & Empreinte factice si le compte n'existe pas &
Le temps de réponse est identique dans les deux cas. \\
Force brute sur la connexion & 5 échecs par compte, 20 par adresse, 15 min &
\cmd{test\_securite.py} \\
Vol de session & \cmd{HttpOnly}, \cmd{SameSite=Strict}, \cmd{Secure}, 2 h &
En-têtes vérifiés par l'analyse dynamique. \\
CSRF & Jeton sur tous les formulaires, actif jusqu'en test &
\cmd{test\_securite.py} \\
Injection de script & CSP \cmd{default-src 'self'}, aucun CDN &
En-tête vérifié par l'analyse dynamique. \\
Vol de la clé de session & Clé du coffre distincte &
\fichier{crypto.py} : deux clés, deux origines. \\
Secret commis dans Git & gitleaks sur tout l'historique &
Travail \cmd{secrets} de la chaîne. \\
Devinage du mot de passe d'une machine & Clé publique, \cmd{PasswordAuthentication no} &
\cmd{Permission denied (publickey)} \\
Évasion de conteneur & \cmd{cap-drop ALL}, \cmd{no-new-privileges} &
\cmd{test\_docker.py} vérifie la commande construite. \\
Épuisement du nœud & \cmd{pids-limit}, \cmd{cpus}, \cmd{memory} &
\cmd{test\_docker.py} \\
Prise de contrôle d'un démon Docker & SSH par clé dédiée, API jamais exposée &
\fichier{workers.py} : \cmd{docker -H ssh://} \\
Faux nœud déclaré & Jeton d'agent exigé &
\cmd{test\_api\_enregistrement\_refuse\_sans\_jeton} \\
Destruction sur un nœud muet & Le silence n'est pas interprété &
\cmd{test\_noeud\_injoignable\_ne\_detruit\_rien} \\
Accès aux machines d'un autre & Filtrage par propriétaire en SQL &
\cmd{test\_securite.py}, cloisonnement entre comptes \\
Collision de port & Index unique partiel en base &
\cmd{test\_contraintes.py} (7 cas) \\
Écoute du trafic & TLS 1.2/1.3, suites ECDHE + AEAD &
Configuration nginx du rôle \cmd{reverse\_proxy}. \\
Dépendance vulnérable & pip-audit, Trivy, SBOM &
Werkzeug 3.1.8 $\rightarrow$ 3.1.9 après alerte. \\
\bottomrule
\end{tabularx}
\end{center}

\begin{retenir}
Dix-huit menaces, dix-huit parades, et pour chacune un endroit où aller voir.
C'est ce dernier point qui compte : une mesure de sécurité qu'aucun test ne
verrouille est une mesure qui disparaîtra à la prochaine modification, sans que
personne s'en aperçoive.
\end{retenir}
